\section{Connections}

\subsection{Connections}

\bdefn

    Let $E\to M$ be a smooth vector bundle over $M$, and denote its space of sections by $\Gamma(E)$.
    A {\emphcolor connection} in $E$ is a map
    $$ \nabla\colon\X(M)\times\Gamma(E)\to\Gamma(E) $$
    written $(X,Y)\mapsto\nabla_XY$ which has the following properties:
    \benum
        \item $\nabla_XY$ is $C^\infty(M)$-linear in $X$:
        $$ \nabla_{f_1X_2+f_2X_2}Y = f_1\nabla_{X_1}Y + f_2\nabla_{X_2},\qquad f_1,f_2\in C^\infty(M). $$
        \item $\nabla_XY$ is $\bR$-linear in $Y$:
        $$ \nabla_X(a_1Y_1+a_2Y_2) = a_1\nabla_XY_1 + a_2\nabla_XY_2,\qquad a_1,a_2\in\bR . $$
        \item $\nabla$ satisfies the product rule: for $f\in C^\infty(M)$:
        $$ \nabla_X(fY) = f\nabla_XY + (Xf)Y . $$
    \eenum

\edefn

\blemm

    If $\nabla$ is a connection over $E$, and $X,\tilde X\in\X(M)$ and $Y,\tilde Y\in\Gamma(E)$ are sections which agree
    on a neighborhood of $p\in M$, then
    $$ \nabla_XY|_p = \nabla_{\tilde X}\tilde Y|_p . $$

\elemm

\bproof

    By linearity it is sufficient to assume $X=0$ or $Y=0$ in a neighborhood of $p$ and show that in each case
    $\nabla_XY=0$.
    For $Y$, suppose $Y$ is zero in a neighborhood $U$ of $p$.
    Let $\psi$ a bump function which is $1$ at $p$ and supported in $U$.
    Thus $\phi Y=0$, and so by the axioms
    $$ 0 = \nabla_X(\phi Y) = \phi\nabla_XY + (X\phi)Y . $$
    Since $Y$ is zero on $\phi$'s support, the second term is zero.
    So evaluating the first term at $p$ gives
    $$ 0 = \phi(p)\nabla_XY|_p = \nabla_XY|_p , $$
    as desired.

    Choosing a similar bump function for $X$ gives
    $$ 0 = \nabla_{\phi X}Y = \phi\nabla_XY $$
    and evaluating at $p$ gives the desired result again.
    \qed

\eproof

\blemm

    Let $\nabla$ be a connection in $E\to M$.
    For an open subset $U\subseteq M$ there is a unique connection $\nabla^U$ on the restricted bundle $E|_U$ such that
    for all $X\in\X(M)$ and $Y\in\Gamma(E)$,
    $$ \nabla_{X|_U}^U(Y|_U) = (\nabla_XY)|_U . $$

\elemm

\bproof

    First we show uniqueness.
    For $X\in\X(U)$ and $Y\in\Gamma(E|_U)$ arbitrary and $p\in U$ we use a bump function to define sections
    $\tilde X\in\X(U)$ and $\tilde Y\in\Gamma(E|_U)$ such that $\tilde X$ and $\tilde Y$ agree with $X,Y$ on a neighborhood
    of $p$ so that
    $$ \nabla_X^UY|_p = \nabla_{\tilde X|_U}^U(\tilde Y|_U)|_p = (\nabla_{\tilde X}\tilde Y)|_p , $$
    showing uniqueness.
    This also demonstrates how to construct $\nabla^U$.
    By the above lemma $\nabla_{\tilde X}\tilde Y|_p$ is independent on choice of $\tilde X,\tilde Y$, and it is easy to see
    that the above formula defines a connection.
    \qed

\eproof

\bnote

    The restricted connection is usually denoted $\nabla$ as well instead of $\nabla^U$.
    This lemma shows there can be no confusion.

\enote

\bprop

    $\nabla_XY$ depends only on $X$'s value at $p$ and $Y$'s values in any neighborhood of $p$.

\eprop

\bproof

    The lemma before this showed that the connection only depends on $X$ and $Y$ in any neighborhood of $p$.
    To show that it depends only on $X_p$, it is sufficient to assume $X_p=0$.
    Let $U$ be a coordinate neighborhood of $p$ and write $X=X^i\ipdv{x^i}$ so that $X^i(p)=0$.
    Then at $p$ by the above lemma on restricting the connection,
    $$ \nabla_XY|_p = \nabla_{X^i\ipdv{x^i}}Y|_p = (X^i\nabla_{\ipdv{x^i}}Y)|_p = 0 $$
    since $X^i(p)=0$.
    \qed

\eproof

Our main concern are connections in the tangent bundle, called just {\emphcolor connections on $M$}.
These are connections
$$ \nabla\colon\X(M)\times\X(M)\to\X(M) . $$

Let $(E_i)$ be a smooth local frame for $TM$ on an open subset $U\subseteq M$.
Then for every $i,j$ we can represent $\nabla_{E_i}E_j$ in $U$ over this coordinate frame.
So there are smooth functions $\Gamma^k_{ij}\in C^\infty(U)$ such that
$$ \nabla_{E_i}E_j = \Gamma^k_{ij}E_k . $$
These are $n^3$ smooth functions ($n=\dim M$) called the {\emphcolor connection coefficients} of $\nabla$ in the given frame.

\bprop

    Let $\nabla$ be a connection on $M$, $E_i$ a local frame over $U$, and $\Gamma^k_{ij}$ be $\nabla$'s connection coefficients
    in this frame.
    Then for $X,Y\in\X(U)$ with representations $X=X^iE_i,Y=Y^jE_j$, one has
    $$ \nabla_XY = (X(Y^k) + X^iY^j\Gamma^k_{ij})E_k . $$

\eprop

\bproof

    Shut up and compute.
    $$ \eqalign{
        \nabla_XY &= \nabla_X(Y^jE_j) = X(Y^j)E_j + Y^j\nabla_XE_j = X(Y^j)E_j + Y^j\nabla_{X^iE_i}E_j\cr
            &= X(Y^j)E_j + X^iY^j\nabla_{E_i}E_j\cr
            &= X(Y^j)E_j + X^iY^j\Gamma^k_{ij}E_k .
    } $$
    \qed

\eproof

\bprop

    Let $\nabla$ be a connection on $M$, and $(E_i)$ and $(\tilde E_j)$ are two local frames over $U$ related by
    $\tilde E_i=A^j_iE_j$ for a matrix $A=(A^j_i)$.
    Let $\Gamma^k_{ij},\tilde\Gamma^k_{ij}$ be the connection coefficients relative to these frames, then
    $$ \tilde\Gamma^k_{ij} = (A^{-1})^k_pA^q_iA^r_j\Gamma^p_{qr} + (A^{-1})^k_pA^q_iE_q(E^p_j) . $$

\eprop

\bproof

    Tedious computations.
    \qed

\eproof

\bprop

    If $TM$ has a global frame $(E_i)$ then any choice of smooth functions $\Gamma^k_{ij}\in C^\infty(M)$ defines a
    connection on $M$ with these as the connection coefficients relative to $(E_i)$.

\eprop

\bproof

    Write $X=X^iE_i$ and $Y=Y^jE_j$, and show that the formula above
    $$ \nabla_XY = (X(Y^k) + X^iY^j\Gamma^k_{ij})E_k $$
    defines a connection on $M$.
    This is clear, or an easy set of computations.
    \qed

\eproof

\bexam

    $T\bR^n$ has a global frame $\ipdv{x^i}$.
    So define $\Gamma^k_{ij}=0$ for all $i,j,k$, this gives us the {\emphcolor Euclidean connection}:
    $$ \bar\nabla_XY = X(Y^i)\pdv{x^i} . $$

\eexam

\bnote

    Linear combinations of connections need not be connections.
    They are obviously still linear, but need not satisfy the product rule.
    For example if $\nabla$ is a connection then
    $$ 2\nabla_X(fY) = 2(XfY + f\nabla_XY) . $$
    But the product rule requires
    $$ (2\nabla_X(fY)) = XfY + f(2\nabla_X)Y = XfY + 2f\nabla_XY , $$
    which are clearly not equal.
    Similarly for sums.

\enote

\bprop

    Every manifold admits a connection on it.

\eprop

\bproof

    Let $M$ be a smooth manifold and cover it with coordinate charts $\set{U_\alpha}$.
    These are all diffeomorphic to Euclidean space and thus can be given the (restriction of the) Euclidean connection
    $\nabla^\alpha$.
    Let $\set{\phi_\alpha}$ be a partition of unity subordinate to $\set{U_\alpha}$ and define
    $$ \nabla_XY = \sum_\alpha\phi_\alpha\nabla^\alpha_XY . $$
    Since partitions of unity are locally finite, this defines a smooth vector field on $M$.
    This is clearly still linear, but we need to show the product rule, as explained in the note above.

    We have
    $$ \nabla_X(fY) = \sum_\alpha\phi_\alpha(XfY + f\nabla^\alpha_XY) = \parens{\sum_\alpha\phi_\alpha}XfY
    + f\sum_\alpha\phi_\alpha\nabla^\alpha_XY = XfY + f\nabla_XY $$
    because the sum of a partition of unity is one.
    \qed

\eproof

\bprop

    Let $\nabla^0,\nabla^1$ be two connections on $M$, then the map $D\colon\X(M)\times\X(M)\to\X(M)$ given by
    $$ D(X,Y) = \nabla^1_XY - \nabla^0_XY $$
    is bilinear over $C^\infty(M)$ and thus defines a $(1,2)$-tensor field called the {\emphcolor difference tensor}.

\eprop

\bproof

    Additivity is clear.
    Multiplicity is by multiplicity in the lower input and the product rule.
    \qed

\eproof

\bthrm

    Let $\nabla^0$ be a connection on $M$, then the set $\A(TM)$ of all connections on $M$ is equal to
    $$ \A(TM) = \set{\nabla^0 + D}[D\in\T^{(1,2)}TM] , $$
    where $\nabla^0+D$ is the connection given by $(\nabla^0+D)_XY=\nabla^0_XY+D(X,Y)$.

\ethrm

\bproof

    The proposition above shows that $\A(TM)$ is contained in this affine space, and it is routine to show that elements of
    the affine space are all connections.
    \qed

\eproof

\subsection{Covariant derivatives of tensor fields}

Recall that a $(1,1)$-tensor over $V$ is just a linear map $V\to V$.
Thus we can naturally define the trace of such a tensor to be the trace of its corresponding linear map.
In general we define the {\emphcolor trace operator}
$$ \tr\colon T^{(k+1,\ell+1)}(V) \to T^{(k,\ell)}(V) $$
by mapping $F\in T^{(k+1,\ell+1)}(V)$ to $\tr(F)\in T^{(k,\ell)}(V)$ where $\tr(F)(\omega^1,\dots,\omega^k,v_1,\dots,v_\ell)$
is defined to be the trace of the $(1,1)$-tensor
$$ F(\omega^1,\dots,\omega^k,\bul,v_1,\dots,v_\ell,\bul) \in T^{(1,1)}(V) . $$
This is clearly a tensor by linearity of the trace.
In terms of a basis
$$ (\tr F)^I_J = F^{I,m}_{J,m} $$
summing over $m$.

Note that we can take the trace of a tensor on different indices, not just the last ones, as long as one is covariant and
the other contravariant.

\bprop

    Let $M$ be a smooth manifold with connection $\nabla$.
    Then $\nabla$ uniquely determines a connection on the tensor bundle $T^{(k,\ell)}TM$ also denoted $\nabla$,
    such that
    \benum
        \item In $T^{(1,0)}TM=TM$, $\nabla$ agrees with the given connection.
        \item In $T^{(0,0)}TM=M\times\bR$ (whose sections are identified with real-valued functions),
        $\nabla$ is ordinary differentiatio of functions,
        $$ \nabla_Xf = Xf . $$
        \item $\nabla$ satisfies the product rule with respect to tensor products,
        $$ \nabla_X(F\otimes G) = (\nabla_XF)\otimes G + F\otimes(\nabla_XG) . $$
        \item $\nabla$ commutes with trace (contractions): if $\tr$ is the trace of a tensor on any pair of indices, one
        covariant and one contravariant, then
        $$ \nabla_X(\tr F) = \tr(\nabla_XF) . $$
    \eenum
    Furthermore the connection satisfies
    \benum
        \item For a vector field $Y$ and covector field $\omega$,
        $$ \nabla_X\iprod{\omega,Y} = \iprod{\nabla_X\omega,Y} + \iprod{\omega,\nabla_XY} $$
        where $\iprod{\omega,Y}$ is just $\omega(Y)$.
        \item For $F\in\T^{(k,\ell)}TM$, covector fields $\omega^1,\dots,\omega^k$, and vector fields $Y^1,\dots,Y^\ell$,
        $$ \eqalign{
            (\nabla_XF)(\omega^1,\dots,\omega^k,Y_1,\dots,Y_\ell) &= X(F(\bar\omega,\bar Y))\cr
            &\hphantom{{}={}}- \sum_{i=1}^k F(\omega^1,\dots,\nabla_X\omega^i,\dots,\omega^k,\bar Y)
             - \sum_{j=1}^\ell F(\bar\omega,Y_1,\dots,\nabla_XY_j,\dots,Y_\ell) .
        } $$
    \eenum

\eprop

\bproof

    First we show that the conditions imply the two properties at the end.
    Note that $\iprod{\omega,Y}=\omega_iY^i=\tr(\omega\otimes Y)$.
    Thus
    $$ \nabla_X\iprod{\omega,Y} = \nabla_X(\tr(\omega\otimes Y)) = \tr\nabla_X(\omega\otimes Y)
    = \tr(\nabla_X\omega\otimes Y + \omega\otimes\nabla_XY) = \iprod{\nabla_X\omega,Y} + \iprod{\omega,\nabla_XY} . $$
    The second property is proven by induction, noting that
    $$ F(\bar\omega,\bar Y) = \tr\cdots\tr(F\otimes\bar\omega\otimes\bar Y) $$
    where we take $k+\ell$ traces, each operating on an upper index of $F$ and the lower index of the corresponding
    covector field, or a lower index of $F$ and the upper index of the corresponding vector field.

    For uniqueness, notice that by the conditions and properties,
    $$ (\nabla_X\omega)(Y) = \iprod{\nabla_X\omega,Y} = \nabla_X(\omega(Y)) - \omega(\nabla_XY) . $$
    Thus we have a unique method of computing the covariant derivatives of covector fields, and by the second property
    this gives us an expression for $\nabla_XF$ for any tensor, showing uniqueness.

    Now to show existence, define $\nabla_X\omega$ for covector fields as in the previous equation, and extend this to
    all tensors.
    Then $\nabla_XF$ is $C^\infty(M)$-multilinear (as can be checked), thus defining a smooth tensor field.
    Checking that $\nabla_X$ actually defines a connection is similarly routine.
    \qed

\eproof

The expressions for the covariant derivative given in the proof/proposition above are not useful for computations.

\bprop

    Let $\nabla$ be a connection on $M$, $(E_i)$ a local frame, $(\epsilon^j)$ its dual coframe, and $\Gamma^k_{ij}$
    the connection coefficients of $\nabla$ with respect to $(E_i)$.
    Let $X=X^iE_i$ be a local vector field.
    \benum
        \item The covariant derivative of $\omega=\omega_i\epsilon^i$ is given by
        $$ \nabla_X\omega = (X(\omega_k)-X^j\omega_i\Gamma^i_{jk})\epsilon^k . $$
    \eenum

\eprop

\TODO{Above.}

\subsection{Vector and tensor fields along maps}

\bdefn

    If $F\colon M\to N$ is a smooth map, and $\pi\colon E\to N$ a vector bundle on $N$, we define its {\emphcolor pullback}
    to be the vector bundle $\pi^*\colon F^*E\to M$ defined by
    $$ F^*E = \set{(p,e)\in M\times E}[F(p)=\pi(e)] \subseteq M\times E , $$
    where $\pi^*$ is projection onto the first component.

\edefn

Projecting onto the second component gives a map $f\colon F^*E\to E$ so that the following diagram commutes

\commsq{F^*E&M\cr E&N}
    {\pi^*}{\pi}{f}{F}

To show that this is indeed a vector bundle we need to give local trivializations.
Let $(U,\Phi)$ be a local trivialization of $E$, where $\Phi\colon\pi^{-1}U\to U\times\bR^k$, then $(F^{-1}U,\Phi^*)$ is a local trivialization of $F^*E$ where
$$ \Phi^*\colon(\pi^*)^{-1}(F^{-1}U) \to F^{-1}U\times\bR^k,\qquad \Phi^*(p,e) = (p,\proj_2(\Phi(e))) . $$
Note that this is well-defined, since $(\pi^*)^{-1}(F^{-1}U)=f^{-1}(\pi^{-1}U)$ so if $(p,e)$ is in it then $e\in\pi^{-1}U$ and $Fp=\pi e$.
Thus $\Phi(e)$ is defined, and $p\in F^{-1}U$.

Now $\Phi^*$ is also bijective.
It is injective since if $\Phi^*(p,e)=\Phi^*(q,e')$ then clearly $p=q$, and so $Fp=\pi e=\pi e'$ so $e,e'\in E_{Fp}$.
Then $\proj_2\Phi|_{E_{Fp}}$ is a bijection $E_{Fp}\to\bR^k$ so $e=e'$.
It is surjective since if $(p,v)\in F^{-1}U\times\bR^k$ then consider the bijection $\Phi|_{E_{Fp}}\colon E_{Fp}\to\bR^k$, there is an $e\in E_{Fp}$ such that
$\Phi(e)=v$.
And $e\in E_{Fp}$ means that $\pi e=Fp$ so $(p,e)$ is the preimage of $(p,v)$.

And $\Phi^*$'s restriction to $(\pi^*)^{-1}\set p=\set{(p,e)}[\pi e=Fp]$ is
$$ \Phi^*(p,e) = (p,\proj_2\Phi(e)) $$
whose second component is $\proj_2\Phi|_{E_{Fp}}$ which is a linear isomorphism to $\bR^k$.

If $\Psi$ is another local trivialization, then
$$ \Phi^*\circ(\Psi^*)^{-1}(p,v) = \Phi^*(p,\Psi^{-1}(Fp,v)) = (p,\proj_2\Phi\Psi^{-1}(Fp,v)) , $$
and the second coordinate is linear, it is equal to $\tau(Fp)v$ where $\tau$ is the transition function between $\Phi,\Psi$.

\bprop

    Sections of the pullback bundle $F^*E$ are equivalent to maps $\sigma\colon M\to E$ where $\sigma(p)\in E_{Fp}$.
    Viewing these as $C^\infty(M)$-modules, they are isomorphic.

\eprop

\bproof

    We want to give a natural correspondence
    $$ \Gamma(F^*E) = \set{\sigma\colon M\to F^*E}[\pi^*\circ\sigma=\id] \oto \set{\sigma\colon M\to E}[\sigma(p)\in E_{Fp}] . $$
    If $\sigma\in\Gamma(F^*E)$ is a section, then define $\tilde\sigma=f\circ\sigma$.
    This is a smooth map and $f\circ\sigma(p)\in E_{Fp}$ since $\pi\circ f\circ\sigma(p)=F\circ\pi^*\circ\sigma(p)=F(p)$.
    Conversely if $\sigma\colon M\to E$ is smooth with $\sigma(p)\in E_{Fp}$ then define $\sigma^*(p)=(p,\sigma(p))$, which is well-defined since $\pi\sigma(p)=Fp$
    by assumption.
    Then $\pi^*\circ\sigma^*(p)=\pi^*(p,\sigma(p))=p$ so $\sigma^*$ is a section.

    These are inverses, since
    $$ (\tilde\sigma)^*(p) = (p,\tilde\sigma(p)) = (p,f\circ\sigma(p)) , $$
    and this is equal to $\sigma(p)$ by definition ($\sigma(p)=(p,e)$ and $e=f(p,e)$).
    And
    $$ \widetilde{\sigma^*}(p) = f\circ\sigma^*(p) = f(p,\sigma(p)) = \sigma(p) . \qed $$

\eproof

Note that this gives us a map $\Gamma(E)\to\Gamma(F^*E)$ mapping $\sigma\in\Gamma(E)$ to $\sigma\circ F$.
In particular if $E_i\colon U\to E$ is a local frame for $E$, then $E_i^*=E_i\circ F$ is a local frame for $F^*E$ in $F^{-1}U$.
This is because $E_i^*|_p=E_i|_{Fp}$ are linearly independent and the rank of $E$ and $F^*E$ are the same.

Now consider connections on pullback bundles, so maps
$$ \nabla\colon \X(M) \times \Gamma(F^*E) \to \Gamma(F^*E) $$
satisfying the conditions of a connection.

\bprop

    Let $F\colon M\to N$ be smooth, $\pi\colon E\to N$ a vector bundle on $N$, and $\pi^*\colon F^*E\to M$ its pullback bundle.
    Further let $\nabla$ be a connection in $E$, then there is a unique connection $\nabla^*$ in $F^*E$ such that
    $$ \nabla_X^*(Y\circ F)|_p = (\nabla_{dF_p(X_p)}Y)|_{Fp} $$
    for all $X\in\X(M)$, $Y\in\Gamma(E)$, and $p\in M$.

\eprop

\bproof

    To show uniqueness, let $p\in M$ and let $(E_i)$ be a local frame of $F(p)\in U\subseteq N$ and consider its pullback frame $(E_i^*)$ of
    $F^*E$.
    Then $\nabla_X^*$ is determined by its values on $E_i^*=E_i\circ F$, which are given by the above formula.
    Thus it is unique, and it is easy to check that $\nabla^*$ is a connection.
    \qed

\eproof

A specific case of interest are pullbacks of the tangent bundle, i.e.\ $F^*TN$.

\bdefn

    Let $F\colon M\to N$ be smooth, then sections of the pullback $F^*TN$ are called {\emphcolor vector fields along $F$}, and the space of sections
    is denoted
    $$ \Gamma(F^*TN) = \X(F) . $$

\edefn

As we saw, a vector field along $F$ amounts to a map $X\colon M\to TN$ such that $X_p\in T_{Fp}N$ for all $p\in M$.

